%------------------------------------------------------------------------------
% Minimal-Template im FOM-Stil, Tectonic-kompatibel
% Bauen:  tectonic -X build
%------------------------------------------------------------------------------
\documentclass[%
	a4paper, 11pt, oneside,
	BCOR=1cm,
	cleardoublepage=plain,
	fleqn,
	numbers=noenddot,
	headings=normal,
	bibliography=totocnumbered,
	listof=numbered,
	ngerman,
	headsepline=0.08em,
]{scrbook}

%--- Basis --------------------------------------------------------------------
\usepackage[TS1,T1]{fontenc}
\usepackage[ngerman]{babel}
\usepackage{csquotes}
\usepackage{amsmath,amsfonts,amssymb}
\usepackage{graphicx}
\usepackage[hyperref,cmyk,dvipsnames]{xcolor}
\usepackage{setspace}
\usepackage{scrlayer-scrpage}

%--- Schriften (der "Look": Times, Helvetica, Bera Mono) ----------------------
\usepackage{mathptmx}
\usepackage[scaled=0.9]{helvet}
\renewcommand{\bfdefault}{b}
\usepackage[scaled=0.9]{beramono}
\usepackage{microtype}
\linespread{1.10}

%--- Ueberschriften und Captions ----------------------------------------------
\setkomafont{sectioning}{\sffamily\bfseries}
\addtokomafont{caption}{\normalfont\normalcolor\small}
\addtokomafont{captionlabel}{\small\bfseries}
\usepackage[format=hang, justification=RaggedRight,
	position=below, width=.9\textwidth]{caption}

%--- Kopf-/Fußzeilen ----------------------------------------------------------
\pagestyle{scrheadings}
\clearpairofpagestyles
\ohead[]{\headmark}
\ofoot[\pagemark]{\pagemark}

%--- Literatur: klassisches BibTeX (in Tectonic eingebaut) ---------------------
% Stil und .bib-Datei werden in _postamble.tex gesetzt.
\usepackage{xurl}

%--- Hyperref -----------------------------------------------------------------
\usepackage[colorlinks=true, linkcolor=black!70, urlcolor=blue,
	bookmarksnumbered=true, breaklinks, plainpages=false]{hyperref}

%--- Metadaten ----------------------------------------------------------------
\newcommand{\daAutor}{Nils Eckerle}
\newcommand{\daTitel}{Ein interessantes Thema}
\newcommand{\daStudiengang}{Bachelor Informatik}
\newcommand{\daGutachter}{Prof. Dr. Gutachter Name}
\newcommand{\daDate}{dd.mm.YYYY}
\newcommand{\daOrt}{Musterstadt}
\newcommand{\daMatrikelnummer}{123456}
\hypersetup{pdftitle={\daTitel}, pdfauthor={\daAutor}}

